% =============================================================================
% Taller 2 - Simulación Montecarlo de un sistema de inventarios
% Contenido del informe. Todas las cifras provienen de General/Resultados.tex,
% generado por tools/exportar_resultados.py (no se digitan a mano).
% =============================================================================

% Histograma de 10 clases a partir de coordenadas (límite inferior, frecuencia)
\newcommand{\histograma}[3]{%
\begin{tikzpicture}
\begin{axis}[
  ybar interval, width=0.95\linewidth, height=5.6cm, ymin=0,
  title={\small #2},
  xlabel={\footnotesize Ganancia anual (millones de COP)},
  ylabel={\footnotesize Frecuencia (réplicas)},
  x tick label style={font=\scriptsize, rotate=40, anchor=north east,
    /pgf/number format/.cd, fixed, precision=1, use comma, 1000 sep={}},
  y tick label style={font=\scriptsize},
  xtick=data, grid=major, grid style={gray!25}]
\edef\histtmp{\noexpand\addplot[fill=#3, draw=black] coordinates {#1};}\histtmp
\end{axis}
\end{tikzpicture}}

\section{Introducción}

Este documento presenta el análisis del Taller~2. Se simula un año completo (los 313 días
hábiles de 2026, de lunes a sábado) de un sistema de inventarios con política $(s,Q)$: al cierre
de cada día, si el inventario final es estrictamente menor que $s$, se piden $Q$ unidades, que
llegan al inicio del siguiente día hábil. La demanda diaria es discreta entre 3 y 7 unidades y
se estudian dos casos: uniforme discreta y empírica \cite{eafit2026}. Cada experimento usa \nReplicas{} réplicas
independientes del año.

El modelo se construyó en dos herramientas que respetan exactamente los mismos supuestos y
ecuaciones del enunciado:

\begin{itemize}
  \item \textbf{Parte A (Excel):} libro con las hojas \texttt{Parametros}, \texttt{Uniforme},
  \texttt{Empirica} y \texttt{Politicas}, con nombres definidos, una fila por día hábil y una tabla
  de datos de 100 réplicas (\texttt{=TABLA(;W2)}).
  \item \textbf{Parte B (Python):} programa \texttt{.py} organizado en funciones (bloques B1 a B8),
  con semilla fija \semilla{} y números aleatorios comunes para comparar políticas.
\end{itemize}

La política recomendada es $\boldsymbol{s=\polS,\ Q=\polQ}$. Las secciones siguientes
justifican esta elección con los resultados de ambas herramientas.

\subsection{Verificación del modelo}

Con $s=6$, $Q=10$ y demandas fijas de 4, 7 y 5 unidades, el libro de Excel (escribiendo
temporalmente esas demandas en \texttt{G6:G8}) y la función \texttt{simulate\_year} de Python
producen exactamente los valores del enunciado (\cref{tab:verificacion}). Además, en ambas hojas
\texttt{B318} muestra 31/12/2026, \texttt{A318} vale 313 y la suma de la columna \texttt{D} es 12
(12 cargos de costo fijo); la comprobación automática de la celda \texttt{U10} es
\comprobacionCalendario.

\begin{table}[H]
\centering
\caption{Ejemplo de verificación ($s=6$, $Q=10$). Idéntico en Excel y en Python.}
\label{tab:verificacion}
\small
\begin{tabular}{cl rrrrrr c r}
\toprule
Día & Fecha & $R$ & $II$ & $D$ & $V$ & $P$ & $IF$ & Pedido & Utilidad del día \\
\midrule
\tablaVerificacionExcel
\bottomrule
\end{tabular}
\end{table}

% =============================================================================
\section{Parte A: análisis del modelo en Excel}
% =============================================================================

\subsection{Construcción del libro}

\begin{itemize}
  \item \texttt{Parametros}: los ocho parámetros en \texttt{B4:B11} con sus nombres definidos
  (\texttt{PuntoReorden}, \texttt{CantPedido}, \texttt{PrecioVenta}, \texttt{CostoUnitario},
  \texttt{CostoPedido}, \texttt{CostoAlm}, \texttt{CostoFijo}, \texttt{AnioSim}) y la tabla
  \texttt{TablaEmpirica} (\texttt{E4:F8}) de la transformada inversa, cuya suma de
  probabilidades se verifica en \texttt{G9}.
  \item \texttt{Uniforme} y \texttt{Empirica}: columnas \texttt{A} a \texttt{R} con las fórmulas
  del enunciado (fila 6 = día 1; fila 7 copiada hasta la 318). La única diferencia es la columna
  \texttt{G}: \texttt{=ALEATORIO.ENTRE(3;7)} frente a
  \texttt{=BUSCARV(ALEATORIO();TablaEmpirica;2;VERDADERO)}, que implementa el método de la transformada inversa \cite{ross2013}.
  \item Resumen de la réplica en \texttt{U6:U8} (más el inventario final del 31/12 en
  \texttt{U9}), tabla de datos \texttt{W6:AA106} con celda de entrada vacía \texttt{W2},
  estadísticos en \texttt{AC5:AC8} y histograma de 10 clases de igual ancho en
  \texttt{AB16:AF26} con su gráfico de columnas (ancho de intervalo 0\,\%).
  \item \texttt{Politicas}: \nPoliticasXl{} combinaciones $(s,Q)$ por cada caso de demanda,
  registradas como valores.
\end{itemize}

\subsection{Pregunta 1: desempeño de la política recomendada}

\Cref{tab:A1} resume las 100 réplicas guardadas en el libro para $s=\polS$, $Q=\polQ$.

\begin{table}[H]
\centering
\caption{Estadísticos de la política recomendada en Excel (100 réplicas).}
\label{tab:A1}
\begin{tabular}{lrr}
\toprule
Estadístico & Demanda uniforme & Demanda empírica \\
\midrule
Ganancia promedio ($\bar G$) & \xlUMedia & \xlEMedia \\
Desviación estándar ($S_G$, \texttt{DESVEST.M}) & \xlUDesv & \xlEDesv \\
Error estándar de la media ($S_G/\sqrt{100}$) & \xlUEE & \xlEEE \\
Probabilidad de ventas perdidas ($\hat p_{\text{perdida}}$) & \xlUPp & \xlEPp \\
Ventas perdidas promedio (unidades/día) & \xlUVp & \xlEVp \\
\bottomrule
\end{tabular}
\end{table}

\begin{figure}[H]
\centering
\begin{subfigure}{0.49\textwidth}
\histograma{\xlUHist}{Excel, demanda uniforme}{EAFIT-blue!55}
\end{subfigure}\hfill
\begin{subfigure}{0.49\textwidth}
\histograma{\xlEHist}{Excel, demanda empírica}{orange!60}
\end{subfigure}
\caption{Histogramas de la ganancia anual en Excel (10 clases de igual ancho entre el mínimo y
el máximo; las frecuencias suman \xlUFreqTotal{} en ambos casos).}
\label{fig:histExcel}
\end{figure}

\textbf{Justificación.} Con $s=\polS$ y $Q=\polQ$ no hubo ningún día con faltantes en ninguna
de las 200 réplicas (313 días cada una), lo que coincide con el resultado teórico de la
pregunta~5. La ganancia anual promedio está alrededor de \$103 millones en ambos casos, con una
desviación estándar cercana a \$3 millones (aproximadamente 3\,\% de la media). El error
estándar de la media, cercano a \$0,3 millones, indica que la ganancia esperada se estima con
una precisión de aproximadamente $\pm$\$0,6 millones (intervalo de confianza del 95\,\%).

\subsection{Pregunta 2: media y varianza teóricas de la demanda}

Para una demanda discreta $D$ con valores $d$ y probabilidades $p(d)$:
$E[D]=\sum_d d\,p(d)$ y $\operatorname{Var}[D]=\sum_d (d-E[D])^2 p(d)$.

\begin{itemize}
  \item \textbf{Uniforme:} $E[D]=0{,}2\,(3+4+5+6+7)=5$ y
  $\operatorname{Var}[D]=0{,}2\,(4+1+0+1+4)=2$ (equivale a $(k^2-1)/12$ con $k=5$ valores);
  desviación estándar $\sqrt2\approx1{,}414$.
  \item \textbf{Empírica:} $E[D]=3(0{,}20)+4(0{,}05)+5(0{,}40)+6(0{,}25)+7(0{,}10)=5$ y
  $\operatorname{Var}[D]=4(0{,}20)+1(0{,}05)+0(0{,}40)+1(0{,}25)+4(0{,}10)=1{,}5$;
  desviación estándar $\approx1{,}225$.
\end{itemize}

Estos valores también se calculan con fórmulas en \texttt{Parametros!A21:D23}.

\textbf{Parecido.} Ambas distribuciones tienen el mismo soporte (3 a 7) y la misma media, 5
unidades por día. Por eso, en promedio, se demandan $5\times313=1565$ unidades al año en ambos
casos y los ingresos esperados son prácticamente iguales.

\textbf{Diferencia.} La empírica tiene 25\,\% menos varianza: concentra el 40\,\% de la
probabilidad en 5 y asigna solo 0,10 a la demanda máxima de 7, frente a 0,20 en la uniforme.

\textbf{Efecto en ventas perdidas y ganancia.} Con $s\ge7$ no hay faltantes en ninguno de los
dos casos (pregunta~5), así que la diferencia no se manifiesta en la política recomendada. Con
$s<7$ sí aparece. Por ejemplo, con $s=6$ solo se pierde venta cuando el día empieza con
exactamente 6 unidades y llega una demanda de 7; como $P(D=7)$ es el doble en la uniforme, la
probabilidad de faltante es aproximadamente el doble (hoja \texttt{Politicas}, \cref{tab:A3}).
En ganancia promedio las diferencias entre casos son pequeñas y del orden del error de
Montecarlo, porque ambas demandas tienen la misma media. La menor varianza de la empírica se
nota en las ventas perdidas, no en el ingreso esperado.

\subsection{Pregunta 3: comparación de políticas y criterio de selección}

\begin{table}[H]
\centering
\caption{Hoja \texttt{Politicas}: \nPoliticasXl{} políticas por caso de demanda, 100 réplicas cada una
(resaltada la recomendada).}
\label{tab:A3}
\small
\begin{tabular}{lcc rr cc}
\toprule
Demanda & $s$ & $Q$ & Ganancia promedio & Desv. estándar & Prob. pérdida & Perdidas/día \\
\midrule
\tablaPoliticasExcel
\bottomrule
\end{tabular}
\end{table}

\textbf{Criterio explícito}, aplicado en orden:
\begin{enumerate}
  \item \textbf{Nivel de servicio:} cero ventas perdidas con ambas demandas. Una venta perdida
  significa un cliente insatisfecho; la pregunta~5 demuestra que $s\ge7$ y $Q\ge7$ lo garantizan
  sin importar la secuencia de demandas.
  \item \textbf{Ganancia:} entre las políticas con servicio completo, la de mayor ganancia
  promedio considerando los dos casos de demanda (promedio de ambos).
  \item \textbf{Riesgo:} si las diferencias de ganancia son menores que el error de Montecarlo,
  se prefiere la de menor desviación estándar y la que respalde la teoría (pregunta~6).
\end{enumerate}

\begin{table}[H]
\centering
\caption{Políticas con cero ventas perdidas en ambos casos, ordenadas por ganancia promedio conjunta.}
\label{tab:A3b}
\small
\begin{tabular}{cc rrr}
\toprule
$s$ & $Q$ & Uniforme & Empírica & Promedio de ambos \\
\midrule
\tablaCriterioExcel
\bottomrule
\end{tabular}
\end{table}

\xlServFrase

Las políticas con $s\le6$ nunca garantizan el servicio. Si solo se mirara la ganancia, la mayor
media de la hoja sería $s=\xlMejorUs$, $Q=\xlMejorUq$ con demanda uniforme (\xlMejorUmedia;
probabilidad de pérdida \xlMejorUpp) y $s=\xlMejorEs$, $Q=\xlMejorEq$ con la empírica
(\xlMejorEmedia; probabilidad de pérdida \xlMejorEpp). En Excel cada política usa números
aleatorios distintos, así que el error estándar de la diferencia de dos medias es aproximadamente
\xlEEdif, y solo diferencias mayores que \xlEEdifDos{} (dos errores estándar) se distinguen del
azar.

\textbf{Costo del servicio completo.} Bajar a $s=6$ sí aumenta un poco la ganancia esperada:
sobra menos inventario al 31/12 (cada unidad sobrante cuesta \$600.000) y una venta perdida solo
cuesta el margen de \$80.000. Con \repsPrecision{} réplicas y números aleatorios comunes
(Parte~B), la ganancia adicional de $s=6$, $Q=\polQ$ frente a la recomendada es \servUSeis{}
$\pm$ \servUSeisIC{} con demanda uniforme y \servESeis{} $\pm$ \servESeisIC{} con la empírica
(diferencias \servUSeisSig{} y \servESeisSig, del \servUSeisPct\,\% y \servESeisPct\,\% de la
ganancia), a cambio de días con faltante en una fracción de \servUSeisPp{} y \servESeisPp{}.
La recomendación $s=\polS$, $Q=\polQ$ paga ese costo pequeño para garantizar servicio completo;
si la empresa aceptara faltantes ocasionales, $s=6$, $Q=\polQ$ sería la alternativa.

\subsection{Pregunta 4: forma de los histogramas y origen de la variabilidad}

Los histogramas de la \cref{fig:histExcel} son unimodales y aproximadamente simétricos, con
forma de campana. La ganancia va de \xlUMin{} a \xlUMax{} millones (uniforme) y de \xlEMin{} a
\xlEMax{} millones (empírica). Es una dispersión moderada: la desviación estándar es de unos
\$3 millones sobre una media de \$103 millones.

La variabilidad entre réplicas se explica sobre todo por el \textbf{inventario final del
31/12/2026}. Por el supuesto~10, ese inventario no tiene valor residual: cada unidad sobrante se
pagó a \$600.000 y nunca se vende. El inventario de cierre depende de en qué punto del ciclo de
pedidos termina el año. Si el último pedido de $Q=\polQ$ unidades llegó en los últimos días,
sobran muchas unidades; si llegó semanas antes, sobran pocas. En las 100 réplicas de Excel varió
entre \xlUInvRango{} (uniforme) y \xlEInvRango{} (empírica) unidades. La relación se midió en la
hoja (celdas \texttt{AC30:AC33}):

\begin{table}[H]
\centering
\caption{Regresión de la ganancia anual sobre el inventario sobrante al 31/12.}
\label{tab:A4}
\begin{tabular}{lrr}
\toprule
 & Uniforme & Empírica \\
\midrule
Inventario final promedio (unidades) & \xlUInvFin & \xlEInvFin \\
Correlación ganancia--inventario final & \xlUCorr & \xlECorr \\
Pendiente (\$ por unidad sobrante) & \xlUPend & \xlEPend \\
Variabilidad explicada ($R^2$) & \xlURR\,\% & \xlERR\,\% \\
\bottomrule
\end{tabular}
\end{table}

La pendiente es del orden del costo unitario de \$600.000: cada unidad que sobra al cierre resta
aproximadamente su costo de compra. Esta sola variable explica el \xlURR\,\% (uniforme) y el
\xlERR\,\% (empírica) de la variabilidad de la ganancia en Excel, y el \pyURR\,\% y \pyERR\,\% en
Python (pendientes de \pyUPend{} y \pyEPend{} por unidad). El resto proviene del volumen total demandado en el año. Su
desviación estándar es $\sqrt{313\operatorname{Var}[D]}\approx25$ unidades (uniforme) y 22
(empírica), y cada unidad vendida deja un margen aproximado de \$80.000 menos el costo de
pedido prorrateado, es decir, cerca de \$2 millones de desviación. La combinación de ambos
efectos ($\sqrt{2^2+2{,}6^2}\approx3{,}3$ millones, tomando el inventario de cierre como
aproximadamente uniforme dentro de un ciclo de $Q$ unidades) reproduce el orden de magnitud de
la desviación estándar observada.

\subsection{Pregunta 5: \texorpdfstring{$s\ge7$ y $Q\ge7$}{s>=7 y Q>=7} implican cero ventas perdidas}

\textbf{Proposición.} Si $s\ge7$ y $Q\ge7$, entonces $P_t=0$ para todo $t=1,\dots,313$ y toda
secuencia de demandas $D_t\in\{3,\dots,7\}$.

\textbf{Demostración.} Basta probar que $II_t\ge7$ para todo $t$. Como $D_t\le7$, eso da
$V_t=\min(II_t,D_t)=D_t$ y por tanto $P_t=D_t-V_t=0$.
\begin{itemize}
  \item \emph{Día 1:} por la ecuación (2), $II_1=R_1=Q\ge7$.
  \item \emph{Día $t\ge2$:} hay dos casos según la revisión del día anterior (ecuación 4).
  \begin{itemize}
    \item Si $IF_{t-1}\ge s$, no se pidió ($O_{t-1}=0$) y $II_t=IF_{t-1}\ge s\ge7$.
    \item Si $IF_{t-1}<s$, se pidió ($O_{t-1}=1$), el pedido llega al inicio del día $t$
    (tiempo de entrega cero) y $II_t=IF_{t-1}+Q\ge0+Q\ge7$, porque $IF_{t-1}\ge0$.
  \end{itemize}
\end{itemize}
En ambos casos $II_t\ge7\ge D_t$. \hfill$\blacksquare$

Las dos condiciones son necesarias. Con $s\le6$ el día puede empezar con 6 unidades sin pedido
pendiente, y una demanda de 7 genera faltante. Con $Q\le6$ un día que empieza con inventario
final 0 más el pedido tiene solo $Q<7$ unidades.

\textbf{Confirmación con los resultados.} En la hoja \texttt{Politicas} todas las políticas con
$s\ge7$ tienen probabilidad de ventas perdidas exactamente 0,0000, y todas las de $s\le6$ tienen
faltantes (\cref{tab:A3}). En Python, las \pyUNCubiertas{} políticas de la rejilla con $s\ge7$ y
$Q\ge7$ tienen cero ventas perdidas con ambas demandas, mientras que las de
$s\le6$ siempre presentan faltantes.

\subsection{Pregunta 6: efecto de aumentar \texorpdfstring{$Q$}{Q} con \texorpdfstring{$s$}{s} fijo}

\begin{table}[H]
\centering
\caption{Excel, $s=\polS$ fijo: ganancia promedio y desviación estándar según $Q$.}
\label{tab:A6}
\begin{tabular}{c rr rr}
\toprule
 & \multicolumn{2}{c}{Ganancia promedio} & \multicolumn{2}{c}{Desviación estándar} \\
\cmidrule(lr){2-3}\cmidrule(lr){4-5}
$Q$ & Uniforme & Empírica & Uniforme & Empírica \\
\midrule
\tablaQExcel
\bottomrule
\end{tabular}
\end{table}

Con $s$ fijo, la ganancia promedio \textbf{primero aumenta y luego disminuye} al crecer $Q$, con
un máximo plano alrededor de $Q\approx13$--17. La desviación estándar crece de forma sostenida.
En Excel la caída para lotes grandes es clara: con $Q=30$ la ganancia es \xlQtreintaU{}
(uniforme) y \xlQtreintaE{} (empírica) menor que con $Q=\polQ$, y la desviación estándar es
\xlQtreintaStdU{} veces mayor. Entre $Q=10$, 15 y 20 las diferencias son del orden del error de
Montecarlo, porque cada política usa números aleatorios distintos. La forma completa de la curva se ve en la rejilla de
Python con $s=7$ y $Q$ de 8 a 30 (\cref{tab:matU,tab:matE}) y en la corrida de \repsPrecision{}
réplicas (\cref{tab:B3}). Con $s\ge7$ no hay ventas perdidas, así que los ingresos y los costos
fijos no dependen de $Q$; la tendencia la explican tres costos:

\begin{itemize}
  \item \textbf{Costo por pedido} (pesa más con $Q$ pequeño): se hacen cerca de $1565/Q$ pedidos
  al año, con un costo de $40.000\times1565/Q$. Ese costo baja de \$6,3 millones con $Q=10$ a
  \$4,2 millones con $Q=15$ y \$3,1 millones con $Q=20$.
  \item \textbf{Inventario sobrante al 31/12} (pesa más con $Q$ grande): en promedio sobran unas
  $Q/2$ unidades adicionales a \$600.000 cada una, es decir, unos \$300.000 por cada unidad
  adicional de $Q$. También aumenta el riesgo, porque el sobrante puede llegar a $s-1+Q$ unidades.
  \item \textbf{Almacenamiento:} \$80 por unidad-día equivale a unos \$25.000 por unidad de
  inventario promedio al año. Es casi despreciable frente a los dos anteriores.
\end{itemize}

Minimizar $C(Q)=40.000\cdot\frac{1565}{Q}+(600.000+25.040)\cdot\frac{Q}{2}$ da
$Q^\ast=\sqrt{2\cdot40.000\cdot1565/625.040}\approx14{,}2$ unidades. Este razonamiento, análogo
al lote económico \cite{winston2004}, coincide con el máximo observado en la simulación y respalda $Q=\polQ$.

% =============================================================================
\section{Parte B: análisis del modelo en Python}
% =============================================================================

El programa reproduce la Parte~A con bucles explícitos: cada línea de \texttt{simulate\_year}
corresponde a una columna de la hoja. Se verifica automáticamente con el ejemplo del enunciado
(\cref{tab:verificacion}) y con el calendario de 313 días. Ejecutar el programa completo (casi
34.000 años simulados) tomó \tiempoPython{} segundos.

\subsection{Pregunta 1: estadísticos, histogramas y comparación con Excel}

\begin{table}[H]
\centering
\caption{Política $s=\polS$, $Q=\polQ$: Python (semilla \semilla) frente a Excel, 100 réplicas cada uno.}
\label{tab:B1}
\small
\begin{tabular}{l rr rr}
\toprule
 & \multicolumn{2}{c}{Demanda uniforme} & \multicolumn{2}{c}{Demanda empírica} \\
\cmidrule(lr){2-3}\cmidrule(lr){4-5}
Estadístico & Python & Excel & Python & Excel \\
\midrule
Ganancia promedio & \pyUMedia & \xlUMedia & \pyEMedia & \xlEMedia \\
Desviación estándar & \pyUDesv & \xlUDesv & \pyEDesv & \xlEDesv \\
Error estándar de la media & \pyUEE & \xlUEE & \pyEEE & \xlEEE \\
Prob. de ventas perdidas & \pyUPp & \xlUPp & \pyEPp & \xlEPp \\
Ventas perdidas por día & \pyUVp & \xlUVp & \pyEVp & \xlEVp \\
\midrule
Diferencia Excel $-$ Python & \multicolumn{2}{c}{\difU{} ($z=$\,\difUZ)} & \multicolumn{2}{c}{\difE{} ($z=$\,\difEZ)} \\
\bottomrule
\end{tabular}
\end{table}

\begin{figure}[H]
\centering
\begin{subfigure}{0.49\textwidth}
\histograma{\pyUHist}{Python, demanda uniforme}{EAFIT-blue!55}
\end{subfigure}\hfill
\begin{subfigure}{0.49\textwidth}
\histograma{\pyEHist}{Python, demanda empírica}{orange!60}
\end{subfigure}
\caption{Histogramas de la ganancia anual en Python (\texttt{bins=10}, mismas clases de igual
ancho que en Excel). El programa los guarda como \texttt{hist\_uniforme.png} y
\texttt{hist\_empirica.png}.}
\label{fig:histPython}
\end{figure}

\textbf{¿Son iguales? ¿Deberían serlo?} No son iguales y \textbf{no deberían serlo}. La lógica
del modelo es idéntica: con las mismas demandas ambas herramientas dan exactamente las mismas
utilidades, como muestra el ejemplo de verificación. Pero cada herramienta genera números
aleatorios distintos (\texttt{ALEATORIO()} en la hoja de cálculo y el generador PCG64 de NumPy
con semilla \semilla{} en Python). Por eso cada una obtiene una muestra diferente de 100 años.

Lo que sí debe ocurrir es que las diferencias sean explicables por el azar. Para dos muestras
independientes, $z=(\bar G_{\text{Excel}}-\bar G_{\text{Python}})/\sqrt{EE_1^2+EE_2^2}$ vale
\difUZ{} y \difEZ, \difZFrase; las medias difieren menos del 1\,\%. Las desviaciones estándar son del mismo orden (unos \$3
millones), las dos herramientas dan cero ventas perdidas y los histogramas tienen la misma forma.

\subsection{Pregunta 2: sensibilidad a la semilla}

\begin{table}[H]
\centering
\caption{Ganancia promedio de 100 réplicas con cinco semillas distintas ($s=\polS$, $Q=\polQ$).}
\label{tab:B2}
\begin{tabular}{r rr}
\toprule
Semilla & Uniforme & Empírica \\
\midrule
\tablaSemillas
\midrule
Rango (máx. $-$ mín.) & \semURango{} (\semURangoPct\,\%) & \semERango{} (\semERangoPct\,\%) \\
Desv. estándar entre semillas & \semUDesv & \semEDesv \\
Error estándar teórico ($S_G/\sqrt{100}$) & \pyUEE & \pyEEE \\
\bottomrule
\end{tabular}
\end{table}

La ganancia promedio cambia entre \semURango{} (uniforme) y \semERango{} (empírica) de una
semilla a otra, como máximo el \semMaxPct\,\% de la media. La desviación estándar entre las cinco medias es
del mismo orden que el error estándar teórico $S_G/\sqrt{n}\approx$ \$0,3 millones, como predice
la teoría: la media de 100 réplicas es una variable aleatoria con esa dispersión.

\textbf{Implicación sobre la precisión.} Con 100 réplicas, la ganancia esperada se conoce con un
intervalo de confianza del 95\,\% de unos $\pm$\$0,6 millones (por ejemplo, \pyUIC{} para la
uniforme). Es suficiente para afirmar que la ganancia anual es de unos \$103 millones, pero
\textbf{no} para distinguir políticas cuyas medias difieren en \$0,1--0,5 millones, como las
mejores de la rejilla. El error disminuye como $1/\sqrt n$ \cite{law2015}: para reducirlo a la mitad hacen falta
4 veces más réplicas. Por eso la Parte~B usa además números aleatorios comunes y una corrida de
\repsPrecision{} réplicas (intervalos de $\pm$\precIC{}).

\subsection{Pregunta 3: rejilla amplia de políticas}

Se evaluaron \nPoliticasPy{} políticas ($s\in\{4,\dots,9\}$ y $Q\in\{8,10,12,\dots,18,20,25,30\}$)
con números aleatorios comunes, 100 réplicas y semilla \semilla. Es una rejilla ocho veces más
amplia que la de Excel (\nPoliticasXl{} políticas). Las \cref{tab:matU,tab:matE} muestran la ganancia promedio en millones de
pesos (negrita: máximo sin restricción; sombreado: política recomendada). La \cref{tab:matPp}
muestra la probabilidad de ventas perdidas.

\begin{table}[H]
\centering
\caption{Python, demanda uniforme: ganancia promedio (millones de COP) por $(s,Q)$.}
\label{tab:matU}
\scriptsize
\setlength{\tabcolsep}{3.5pt}
\begin{tabular}{c|*{12}{r}}
\toprule
$s\backslash Q$ & \matUCabeza \\
\midrule
\matUMedia
\bottomrule
\end{tabular}
\end{table}

\begin{table}[H]
\centering
\caption{Python, demanda empírica: ganancia promedio (millones de COP) por $(s,Q)$.}
\label{tab:matE}
\scriptsize
\setlength{\tabcolsep}{3.5pt}
\begin{tabular}{c|*{12}{r}}
\toprule
$s\backslash Q$ & \matECabeza \\
\midrule
\matEMedia
\bottomrule
\end{tabular}
\end{table}

\begin{table}[H]
\centering
\caption{Python: probabilidad de ventas perdidas por $(s,Q)$ (uniforme arriba, empírica abajo).}
\label{tab:matPp}
\scriptsize
\setlength{\tabcolsep}{3.5pt}
\begin{tabular}{c|*{12}{r}}
\toprule
$s\backslash Q$ & \matUCabeza \\
\midrule
\matUPp
\midrule
\matEPp
\bottomrule
\end{tabular}
\end{table}

\textbf{¿Coincide con Excel?} Por ganancia pura, la mejor política de la rejilla es $s=\pyUMejorS$,
$Q=\pyUMejorQ$ con demanda uniforme (\pyUMejorMedia) y $s=\pyEMejorS$, $Q=\pyEMejorQ$ con la
empírica (\pyEMejorMedia). Las dos tienen $s\le6$ y pierden ventas (probabilidad \pyUMejorPp{} y
\pyEMejorPp{}), así que no cumplen el criterio de servicio. Su ventaja sobre la política
recomendada, medida con diferencias pareadas, es \pyUDifPar{} $\pm$ \pyUDifParIC{}
(\pyUDifParSig) y \pyEDifPar{} $\pm$ \pyEDifParIC{} (\pyEDifParSig; IC del 95\,\%). Es el mismo
efecto que mide la corrida de \repsPrecision{} réplicas de la Parte~A (pregunta~3): un $s$ menor
gana algo más porque deja menos inventario sin valor al 31/12, pero pierde ventas.

Con el mismo criterio de la Parte~A (cero ventas perdidas y máxima ganancia promedio conjunta),
la rejilla de 100 réplicas elige $s=\pyServS$, $Q=\pyServQ$. Esta política difiere de la
recomendada solo en $Q$, y por una cantidad que no es significativa (\pyUDifServ{} $\pm$
\pyUDifServIC{} y \pyEDifServ{} $\pm$ \pyEDifServIC). Para resolver el empate se repitieron las
políticas con $s=7$ y $Q=11,\dots,19$ con \repsPrecision{} réplicas y números aleatorios comunes
(\cref{tab:B3}).

\begin{table}[H]
\centering
\caption{$s=7$ con \repsPrecision{} réplicas: ganancia promedio y diferencia pareada frente a
$Q=\polQ$ (IC del 95\,\%).}
\label{tab:B3}
\scriptsize
\begin{tabular}{c rr rr}
\toprule
 & \multicolumn{2}{c}{Demanda uniforme} & \multicolumn{2}{c}{Demanda empírica} \\
\cmidrule(lr){2-3}\cmidrule(lr){4-5}
$Q$ & Media & Dif. vs.\ $Q=\polQ$ & Media & Dif. vs.\ $Q=\polQ$ \\
\midrule
\tablaPrecision
\bottomrule
\end{tabular}
\end{table}

Con mayor precisión se confirma una meseta: \precMejoresFrase. Las medias bajan hacia los
extremos, y la caída ya es significativa en \precPeores. Esto coincide con el análisis de costos de la
pregunta~6 de la Parte~A ($Q^\ast\approx14$). Por lo tanto, entre las políticas con servicio completo,
\textbf{la mejor coincide con la elegida en Excel}: $s=\polS$, $Q=\polQ$. Es el menor $s$ que garantiza servicio completo
y un $Q$ en el centro de la meseta óptima. Las diferencias de orden en la rejilla de 100
réplicas (la recomendada queda en el puesto \pyURango{} con demanda uniforme y \pyERango{} con
la empírica) se deben al ruido de Montecarlo y no a diferencias reales.

\subsection{Pregunta 4: ¿por qué la misma semilla para todas las políticas?}

\texttt{compare\_policies} crea el generador con la misma semilla para cada política
(\emph{números aleatorios comunes} \cite{law2015,ross2013}). Así, todas las políticas enfrentan exactamente las mismas
100 secuencias de 313 demandas, y la diferencia de ganancia entre dos políticas se debe solo a la
política, no a que una tuvo por suerte un año de demanda más alta. Estadísticamente,
\[
\operatorname{Var}(\bar G_A-\bar G_B)=\operatorname{Var}(\bar G_A)+\operatorname{Var}(\bar G_B)-2\operatorname{Cov}(\bar G_A,\bar G_B),
\]
y con números comunes la covarianza es positiva (años de alta demanda favorecen a ambas
políticas), lo que reduce la varianza de la diferencia.

\textbf{¿Qué pasaría sin ellos?} Las diferencias mezclarían el efecto de la política con el
ruido del muestreo, y el ranking podría invertirse por azar. Por ejemplo, al comparar la mejor
política sin restricción con la recomendada usando corrientes aleatorias independientes, el
intervalo de la diferencia fue $\pm$\pyUDifIndIC{} (uniforme) y $\pm$\pyEDifIndIC{}
(empírica), frente a $\pm$\pyUDifParIC{} y $\pm$\pyEDifParIC{} con números comunes. La
diferencia estimada también cambió (\pyUDifInd{} y \pyEDifInd{}, frente a \pyUDifPar{} y
\pyEDifPar{}). Con 72 políticas evaluadas de forma independiente, la que parezca mejor sería, en
buena medida, la que tuvo más suerte. Los números comunes no eliminan el error de Montecarlo,
pero hacen la comparación justa y más precisa con el mismo número de réplicas.

\subsection{Pregunta 5: Excel frente a Python}

\begin{table}[H]
\centering
\caption{Comparación de las dos herramientas para este modelo.}
\label{tab:B5}
\small
\begin{tabularx}{\textwidth}{>{\raggedright\arraybackslash}p{2.6cm}>{\raggedright\arraybackslash}X>{\raggedright\arraybackslash}X}
\toprule
Criterio & Excel & Python \\
\midrule
Tiempo de cómputo &
La tabla de datos recalcula el libro completo (2 hojas de 313 filas y 18 columnas) una vez por
réplica: 100 recálculos por cada F9. El libro se vuelve lento, por eso se usa el modo
``automático excepto en las tablas de datos''. &
El programa completo (unos 34.000 años simulados: estadísticos, 5 semillas, 72 políticas
$\times$ 2 demandas y la corrida de \repsPrecision{} réplicas) tarda \tiempoPython{} s. \\
\addlinespace
Cambiar réplicas o políticas &
Más réplicas exige ampliar la tabla de datos y los rangos de las fórmulas. Cada política exige
cambiar $s$ y $Q$ a mano, recalcular y copiar valores (\nPoliticasXl{} políticas = \nPoliticasXl{} ciclos manuales). &
Se cambia \texttt{N\_REPLICATIONS} o la lista \texttt{POLICIES} en una línea. Pasar de \nPoliticasXl{} a \nPoliticasPy{}
políticas no requiere trabajo adicional. \\
\addlinespace
Trazabilidad &
Excelente para ver un año: cada día y cada costo están en una celda visible y auditable. Pero las
réplicas de la tabla de datos no se pueden reproducir: \texttt{ALEATORIO()} no tiene semilla y
cambia con cada F9. &
Cada resultado se puede reproducir exactamente gracias a la semilla. El flujo
(calendario $\to$ demanda $\to$ simulación $\to$ estadísticos) es explícito y puede imprimir el
registro diario (\texttt{keep\_log=True}). \\
\addlinespace
Riesgo de errores &
Referencias mal copiadas, rangos que no cubren las 313 filas, celda de entrada de la tabla de datos
mal elegida, valores pegados que quedan desactualizados. &
Errores de lógica (índices, orden de eventos). Se mitigan con funciones pequeñas, la prueba
obligatoria automática (\texttt{check\_statement\_example}) y la verificación del calendario. \\
\addlinespace
Comparación de políticas &
Sin números aleatorios comunes: cada política usa números distintos, lo que añade ruido. &
Números aleatorios comunes, intervalos pareados y miles de réplicas sin esfuerzo adicional. \\
\bottomrule
\end{tabularx}
\end{table}

En síntesis, Excel es ideal para construir, explicar y auditar la lógica de un año. Python es
superior para la experimentación: muchas réplicas y políticas, resultados reproducibles y
comparaciones estadísticamente rigurosas. Validar un modelo contra el otro, con el ejemplo de
verificación y con estadísticos comparables, aumenta la confianza en ambos.

\section{Conclusiones}

\begin{enumerate}
  \item Los dos modelos implementan las mismas ecuaciones y reproducen exactamente el ejemplo de
  verificación del enunciado.
  \item Se recomienda la política $\boldsymbol{s=\polS,\ Q=\polQ}$. Garantiza cero ventas
  perdidas con cualquier secuencia de demandas (demostrado en la pregunta~5 y confirmado en las
  simulaciones), y entre las políticas con servicio completo su ganancia, de unos \$103 millones
  al año, es estadísticamente equivalente a la mejor de una rejilla de \nPoliticasPy{} políticas.
  Las políticas con $s=6$ ganan cerca de \servGananciaMax{} más al año (alrededor del 0,4\,\%) a
  costa de perder ventas.
  \item La variabilidad de la ganancia anual (desviación estándar de unos \$3 millones) se debe
  sobre todo al inventario sin valor residual que queda al 31 de diciembre, y en segundo lugar al
  volumen total de demanda del año.
  \item Las dos distribuciones de demanda tienen la misma media (5) y producen ganancias
  promedio similares. La empírica, con menor varianza (1,5 frente a 2), genera menos ventas
  perdidas cuando $s<7$.
  \item Con 100 réplicas la ganancia esperada se estima con un margen de unos $\pm$\$0,6
  millones (95\,\%). Para comparar políticas cercanas se necesitan números aleatorios comunes y
  más réplicas.
\end{enumerate}

